\chapter{人力的设计：移民、劳工与准证体系}

\section{两个世界，一座城市}

在新加坡停留的这些天里，有一个现象始终萦绕在我心头。

在街头、工地，以及各类建筑工程的现场，我看到大量来自孟加拉国和印度的劳工，他们身穿统一的工装，成群结队地出现在城市的各个角落。而在会议场合与工作场合，我接触到的则是另一类人——来自中国大陆的半导体工程师，就职于全球知名的芯片公司，收入体面，生活稳定。

这两类人共享同一座城市，却运行在几乎完全平行的轨道上，彼此之间几乎没有任何交集。

这并非偶然，而是制度设计的结果。

\section{三层准证体系}

新加坡对外来劳动力的管理，依托一套精密的工作准证体系，大致分为三个层级。

\subsection*{第一层：Work Permit（工作准证）}

这一类准证主要针对从事建筑、清洁、船厂等体力劳动的工人。持证者的月薪通常在 900 至 1,500 新加坡元之间，住宿由雇主统一提供，不得携带家属，也几乎没有向上流动的制度通道。

现实情况往往更为严峻——许多工人在出发前便已背负高额中介费，抵达后面对的是漫长的工时和有限的权利保障。

\subsection*{第二层：S Pass（S 准证）}

这一层级面向中等技术水平的从业者，例如技术员和操作员。月薪区间约为 3,000 至 5,000 新加坡元，持证者可以自行租房居住，但数量受到行业配额的限制。

\subsection*{第三层：Employment Pass（就业准证）}

这是面向高端人才的准证类别，涵盖工程师、科研人员及各类专业人士。月薪通常在 5,000 新加坡元以上，高端岗位甚至可达 15,000 元以上。持证者可携带家属，并具备申请永久居民（PR）乃至公民身份的资格。

\section{芽笼：系统的缓冲地带}

在新加坡，有一个地名在讨论外劳议题时频繁出现，那便是芽笼（Geylang）。

这一区域廉价住房密集，交通便利，围绕着外来劳动者的日常需求形成了一套完整的生活配套。每逢周日，当大多数工地停工休息，大批外劳便聚集于此，形成颇为壮观的人流景象。芽笼在新加坡城市规划中扮演着一种"缓冲区"的角色——它在空间上将底层劳动力的生活圈与城市的精英区域隔离开来，维持着整座城市表面上的有序与整洁。

\section{住房的分层逻辑}

外来劳动者的居住安排，与其准证级别高度对应。

持 Work Permit 的工人通常居住在由雇主提供的集体宿舍或廉价出租屋，生活空间逼仄，自主性有限。而持 S Pass 或 Employment Pass 的人员，则可以租住政府组屋（HDB）或私人公寓，融入更为正常的城市生活。

住在哪里，由手中的那张准证决定。

\section{成本的分配}

在这套体系中，企业与工人各自承担不同的成本。

雇主方面，除支付工资外，还需负担工人的住宿、基本医疗，以及向政府缴纳的外劳税（levy）。这笔税款的设计，既是对企业雇用低薪外劳的一种调节机制，也是政府财政的一项来源。

工人方面，则需自行承担出发前的中介费用，以及在新加坡期间的日常生活开销。对于底层劳工而言，这两项支出往往构成沉重的经济压力。

\section{一个高度工程化的人力系统}

综合来看，新加坡构建的是一套高度工程化的人力资源系统。

在低端，它以成本控制为核心，通过 Work Permit 机制引入廉价劳动力，维系城市的基础建设与服务运转；在中端，它以结构平衡为目标，通过 S Pass 的配额制度调节中层技术人员的供给；在高端，它以人才吸引为导向，以就业准证和移民通道为工具，将全球各地的专业人才纳入本国的经济体系之中。

这套系统的高效运转，建立在严格的分层管理与精密的制度设计之上。它使新加坡得以在土地和人口极为有限的条件下，同时维持城市基础设施的持续更新与高端产业的持续发展。

然而，这套系统的另一面，是数以万计的底层劳动者，在这座光鲜的城市中，过着与其光鲜程度极不相称的生活。这座城市的高效，并非没有代价——只是这代价，由谁来承担，又由谁来看见，本身便是一个值得追问的问题。
